\documentclass[12pt,a4paper]{article}

\usepackage[utf8]{inputenc}
\usepackage[T1]{fontenc}
\usepackage[french]{babel}
\usepackage{geometry}
\geometry{margin=2.5cm}
\usepackage{hyperref}
\hypersetup{
    colorlinks=true,
    linkcolor=blue,
    urlcolor=cyan,
    citecolor=blue,
}
\usepackage{enumitem}
\usepackage{booktabs}
\usepackage{xcolor}
\usepackage{titlesec}
\usepackage{fancyhdr}
\usepackage{float}
\usepackage{tabularx}

\pagestyle{fancy}
\fancyhf{}
\lhead{Rapport d'Avancement n°4 - SpeechCoach}
\rhead{Master SIE - 2026}
\cfoot{\thepage}

\titleformat{\section}{\large\bfseries\color{blue!70!black}}{}{0em}{}[\titlerule]
\titleformat{\subsection}{\bfseries\color{blue!50!black}}{}{0em}{}

\title{
    \vspace{-2cm}
    \large Master Systèmes Intelligents pour l'Éducation (SIE)\\
    \vspace{0.5cm}
    \huge \textbf{Rapport d'État d'Avancement n°4}\\
    \large \textit{Itérations sur le bloc d'exercices, la couche ASR et l'analyse visuelle device-aware}
}
\author{\textbf{Étudiant :} Hicham Moussaid \\ \textbf{Encadrant :} Pr. Abdelaaziz Hessane}
\date{Période : après la dernière réunion de suivi}

\begin{document}

\maketitle

\section{Objectif de la phase}
Cette phase avait deux objectifs :
\begin{itemize}
    \item améliorer le bloc d'exercices pour qu'il soit plus crédible et plus utile dans le rapport final ;
    \item améliorer la détection des hésitations et des répétitions dans le pipeline audio, en particulier dans les cas français, anglais et arabe.
\end{itemize}

\section{Travail sur le bloc d'exercices}

\subsection{Problème initial}
La première version du bloc d'exercices présentait plusieurs limites :
\begin{itemize}
    \item le plan sur 7 jours était trop long ;
    \item le bloc apparaissait presque systématiquement, même pour des bonnes sessions ;
    \item certaines consignes répétaient le bilan ou la priorité au lieu d'apporter une vraie action ;
    \item des problèmes de setup comme la netteté ou le cadrage étaient parfois affichés comme des exercices de prise de parole.
\end{itemize}

\subsection{Itérations réalisées}
Le travail n'a pas été fait en une seule étape. Plusieurs ajustements successifs ont été nécessaires :
\begin{itemize}
    \item remplacement du plan par défaut sur 7 jours par une logique conditionnelle ;
    \item introduction de plusieurs modes : \texttt{none}, \texttt{light\_tip}, \texttt{setup\_action}, \texttt{single\_exercise}, \texttt{mini\_plan\_3\_days} ;
    \item réduction du plan long vers un mini-plan sur 3 jours ;
    \item séparation entre \textbf{exercice de parole} et \textbf{action de préparation} ;
    \item suppression de nombreuses répétitions textuelles entre exercice, bilan, priorité et encouragement ;
    \item réécriture de plusieurs formulations jugées trop scolaires, trop vagues ou trop artificielles ;
    \item ajustement du comportement pour que les sessions fortes n'affichent plus automatiquement un exercice.
\end{itemize}

\subsection{Résultat retenu}
Le comportement obtenu est plus cohérent :
\begin{itemize}
    \item les très bonnes sessions peuvent ne plus afficher d'exercice ;
    \item les cas intermédiaires affichent un conseil court ou un exercice ciblé ;
    \item les cas plus faibles affichent un mini-plan plus lisible ;
    \item les problèmes techniques sont traités comme des vérifications et non comme des exercices oraux.
\end{itemize}

\section{Travail sur l'ASR : hésitations et répétitions}

\subsection{Constat de départ}
Le problème principal observé était le suivant : la métrique \texttt{Hesitations} ne reflétait pas toujours ce qui était réellement prononcé. Dans plusieurs tests, des fillers comme \textit{um}, \textit{uh}, \textit{euh}, \textit{hum}, \textit{mmm} ou certaines formes arabes étaient entendus, mais pas correctement conservés dans la transcription \cite{radford2022robust,ctranslate2}. Cela affectait directement le score voix/débit.

\subsection{Première piste explorée : changer de moteur}
Un benchmark isolé a été mené avec \textbf{CrisperWhisper} \cite{zusag2024crisperwhisper,crisperwhisper_repo}. L'idée était de vérifier si un moteur plus orienté verbatim pouvait résoudre le problème à la racine.

\subsection{Résultat du benchmark CrisperWhisper}
La piste a été rejetée, car :
\begin{itemize}
    \item le temps d'exécution était beaucoup trop élevé ;
    \item l'amélioration sur nos cas réels n'était pas convaincante ;
    \item la qualité globale des sorties n'était pas suffisamment stable.
\end{itemize}

La conclusion a été de conserver Faster-Whisper et de corriger notre propre pipeline.

\subsection{Problèmes réellement identifiés}
L'analyse détaillée a montré que le problème venait de plusieurs sources :
\begin{itemize}
    \item certaines chaînes arabes dans les prompts et les regex étaient corrompues ;
    \item le prompt multilingue de secours favorisait les hallucinations ;
    \item certaines listes de fillers étaient incomplètes ou mal calibrées ;
    \item Faster-Whisper lissait encore une partie des sons nasaux de type \textit{mm}, \textit{mmm}, \textit{همم} ;
    \item les répétitions simples et les redémarrages n'étaient pas assez bien modélisés.
\end{itemize}

\subsection{Itérations réalisées sur Faster-Whisper}
Le correctif final est le résultat de plusieurs itérations :
\begin{itemize}
    \item restauration des chaînes arabes correctes dans les prompts et les regex ;
    \item suppression du prompt multilingue global ;
    \item mise en place d'une stratégie en deux passes en mode auto :
    \begin{itemize}
        \item première passe pour détecter la langue ;
        \item seconde passe avec langue forcée et prompt spécifique à la langue détectée ;
    \end{itemize}
    \item amélioration des listes de fillers en anglais, français et arabe ;
    \item exclusion de certains marqueurs discursifs trop généraux, par exemple \textit{like} et \textit{you know} en anglais, ou \textit{donc}, \textit{du coup}, \textit{en fait}, \textit{enfin} en français ;
    \item conservation de certains fillers jugés pédagogiquement pertinents, comme \textit{ben} en français ;
    \item ajout d'un détecteur acoustique léger pour les fillers nasaux de type \textit{hum}, \textit{hmm}, \textit{mm}, \textit{mmm}, \textit{مم}, \textit{همم}, à partir de caractéristiques extraites avec \texttt{librosa} \cite{mcfee2015librosa} ;
    \item amélioration du comptage des répétitions via un détecteur spécifique de redémarrages et répétitions simples.
\end{itemize}

\subsection{Résultats observés par langue}
\textbf{Anglais}
\begin{itemize}
    \item l'anglais est devenu le cas le plus stable ;
    \item les fillers comme \textit{um}, \textit{uh}, \textit{mmm}, \textit{hmm} sont mieux conservés ;
    \item les répétitions comme \textit{like like like} ou \textit{the system, the system} sont correctement détectées.
\end{itemize}

\textbf{Français}
\begin{itemize}
    \item plusieurs itérations ont été nécessaires ;
    \item un bug dans les regex françaises a d'abord empêché le bon comptage ;
    \item les fillers de type \textit{hum}, \textit{hmm}, \textit{mmm} étaient parfois perdus par la transcription ;
    \item le fallback acoustique a amélioré la situation ;
    \item le comptage final est plus cohérent après nettoyage de la liste des marqueurs à conserver.
\end{itemize}

\textbf{Arabe}
\begin{itemize}
    \item l'arabe restait le cas le plus difficile ;
    \item les fillers textuels comme \textit{يعني}, \textit{آه}, \textit{ام} sont mieux récupérés ;
    \item les formes nasales comme \textit{مم} ou \textit{همم} ne sont pas toujours transcrites ;
    \item le détecteur acoustique permet néanmoins de récupérer une partie de ces cas même lorsqu'ils n'apparaissent pas dans le texte.
\end{itemize}

\subsection{Impact sur la note voix/débit}
Un point important mis en évidence pendant cette phase est que la note voix/débit dépend directement de la configuration ASR. Dès qu'on change le modèle ou le comportement de transcription, on modifie :
\begin{itemize}
    \item le nombre d'hésitations détectées ;
    \item le nombre de répétitions détectées ;
    \item parfois même certaines métriques temporelles dérivées.
\end{itemize}

L'ASR a donc un impact direct sur l'évaluation pédagogique globale.

\section{Comparaison de configurations Faster-Whisper}

\subsection{Configurations comparées}
Deux configurations ont été testées :
\begin{itemize}
    \item \texttt{base + beam size 5} ;
    \item \texttt{medium + beam size 5}.
\end{itemize}

\subsection{Résultat}
La configuration \texttt{base + beam size 5} est plus rapide, mais elle dégrade trop la fidélité utile pour SpeechCoach :
\begin{itemize}
    \item perte d'une partie des fillers en anglais ;
    \item résultats plus faibles en arabe ;
    \item transcriptions globalement moins fiables pour l'évaluation pédagogique ;
    \item variation artificielle du score voix/débit.
\end{itemize}

À l'inverse, \texttt{medium + beam size 5} est plus coûteuse, mais elle stabilise mieux le comptage des hésitations et des répétitions. C'est donc la configuration la plus adaptée à l'objectif du projet.

\section{Travail sur le regard caméra selon le type d'appareil}

\subsection{Problème observé}
Après stabilisation du bloc d'exercices et de la couche ASR, le point le plus fragile restait l'analyse visuelle, en particulier le regard caméra. Une même logique générique était appliquée à tous les cas, alors que les conditions de capture changent fortement entre :
\begin{itemize}
    \item un ordinateur portable ou une webcam intégrée ;
    \item un smartphone en mode portrait ;
    \item une tablette utilisée en mode selfie.
\end{itemize}

Les premiers essais ont montré que cette hypothèse unique pénalisait surtout les vidéos smartphone. Dans un cas de test où le visage était centré, les mains visibles et le regard dirigé vers la caméra, le système produisait pourtant des valeurs très faibles pour la présence, le contact visuel et la visibilité des mains. Le problème ne venait donc pas seulement du seuil regard, mais plus globalement du pipeline vision appliqué aux vidéos mobiles.

\subsection{Première version : profils par appareil}
Une première amélioration a consisté à introduire un contexte d'appareil dans le pipeline :
\begin{itemize}
    \item \texttt{laptop\_desktop} ;
    \item \texttt{tablet} ;
    \item \texttt{smartphone} ;
    \item \texttt{unknown}.
\end{itemize}

Le choix a été propagé depuis le studio jusqu'au traitement backend afin que l'analyse visuelle puisse adapter ses seuils. Le principe retenu est volontairement simple :
\begin{itemize}
    \item \texttt{laptop\_desktop} conserve un comportement proche du réglage initial ;
    \item \texttt{tablet} tolère davantage certains décalages latéraux ;
    \item \texttt{smartphone} est traité à part, car c'est le cas où la géométrie et le cadrage diffèrent le plus.
\end{itemize}

\subsection{Correctifs spécifiques smartphone}
Le travail sur smartphone a montré que le problème n'était pas uniquement un problème de seuils. Trois ajustements ont été nécessaires :
\begin{itemize}
    \item conservation du format portrait lors de l'extraction des frames, au lieu de forcer un redimensionnement trop proche d'un format paysage ;
    \item augmentation de la fréquence d'échantillonnage des frames pour smartphone ;
    \item assouplissement des tolérances de regard pour ce mode.
\end{itemize}

Ces ajustements ont permis de faire évoluer un cas de test smartphone d'un résultat incohérent à un résultat crédible. Après correctif, le même clip a produit :
\begin{itemize}
    \item \textbf{Présence visage} : de 71\% vers 100\% ;
    \item \textbf{Contact visuel} : de 10\% vers plus de 80\% ;
    \item \textbf{Mains visibles} : de 7\% vers plus de 60\%.
\end{itemize}

\subsection{Résultats sur laptop et tablette}
Les tests laptop et tablette se sont révélés plus stables que le smartphone :
\begin{itemize}
    \item sur laptop, les résultats explicites en mode \texttt{laptop\_desktop} étaient déjà cohérents et ont confirmé que le réglage initial était proche du bon comportement ;
    \item sur tablette, le regard et la présence ont également donné des résultats crédibles, même dans un cadrage selfie où les mains n'étaient presque jamais visibles ;
    \item l'interprétation posture/gestes reste néanmoins à affiner plus tard pour éviter de sanctionner trop fortement les plans serrés où les mains sont naturellement hors cadre.
\end{itemize}

Une dernière passe d'ajustement a ensuite été appliquée pour mieux fermer cette phase :
\begin{itemize}
    \item assouplissement supplémentaire des seuils tablette afin de mieux couvrir les usages réels en paysage ;
    \item ajout d'une heuristique géométrique finale permettant de faire basculer certains formats paysage standards vers le profil \texttt{laptop\_desktop} lorsque les indices tablette ou smartphone sont absents.
\end{itemize}

\subsection{Mode \texttt{unknown}}
Une seconde itération a consisté à rendre le mode \texttt{unknown} plus intelligent, afin de ne pas dépendre uniquement d'une sélection manuelle côté utilisateur.

La stratégie retenue est hiérarchique :
\begin{itemize}
    \item lecture des métadonnées vidéo et des tags via \texttt{ffprobe} ;
    \item prise en compte de la rotation réelle de la vidéo ;
    \item utilisation du nom d'origine du fichier comme indice supplémentaire ;
    \item recours à une heuristique géométrique simple pour les cas paysage les plus standards ;
    \item repli sur le profil \texttt{unknown} si aucun indice n'est suffisamment fiable.
\end{itemize}

Les résultats finaux obtenus montrent que :
\begin{itemize}
    \item le smartphone peut désormais être auto-détecté via les métadonnées et la rotation ;
    \item la tablette peut également être reconnue lorsqu'un indice matériel est disponible dans les métadonnées ou dans le nom du fichier ;
    \item le laptop/desktop peut être déduit proprement par heuristique dans les formats paysage standards lorsqu'aucun marqueur smartphone ou tablette n'est présent ;
    \item les tests menés en fin de phase ont permis de valider un cas fonctionnel dans les trois situations : smartphone, tablette et laptop.
\end{itemize}

En pratique, cela signifie que le mode \texttt{unknown} n'est plus uniquement un profil de secours : il devient un mode de détection automatique exploitable dans les cas les plus courants.

Plus explicitement, si l'appareil est détecté correctement, le pipeline bascule directement vers le profil déjà défini correspondant : \texttt{smartphone}, \texttt{tablet} ou \texttt{laptop\_desktop}. Si aucun indice n'est jugé suffisamment fiable, le système ne force pas une classe ; il conserve alors le profil \texttt{unknown}, qui joue le rôle de fallback générique avec des seuils neutres et un traitement non spécialisé.

\section{Synthèse}
\begin{table}[H]
    \centering
    \begin{tabularx}{\textwidth}{@{}p{5.4cm}p{2.5cm}X@{}}
        \toprule
        \textbf{Sujet} & \textbf{État} & \textbf{Observation} \\
        \midrule
        Bloc d'exercices adaptatif & Implémenté & Plus court, conditionnel et mieux aligné avec le niveau réel de la session. \\
        Plan 7 jours & Abandonné & Remplacé par des formats plus compacts et plus crédibles. \\
        Benchmark CrisperWhisper & Rejeté & Trop lent et peu convaincant sur nos cas réels. \\
        Hésitations multilingues & Améliorées & Travail itératif sur prompts, regex et fallback acoustique. \\
        Répétitions & Améliorées & Détection plus robuste des répétitions simples et redémarrages. \\
        Configuration ASR retenue & Validée & Retour à \texttt{medium + beam size 5}. \\
        Regard caméra par appareil & Implémenté & Profils distincts pour laptop, smartphone et tablette. \\
        Mode \texttt{unknown} & Validé & Auto-détection fonctionnelle par métadonnées, rotation, nom du fichier et heuristique paysage. \\
        \bottomrule
    \end{tabularx}
    \caption{Résumé des travaux de cette phase.}
\end{table}

\section{Prochaines étapes}
\begin{itemize}
    \item consolider l'interprétation posture/gestes sur les plans serrés smartphone et tablette ;
    \item continuer à tester le mode \texttt{unknown} sur davantage de cas réels ;
    \item continuer à documenter les résultats backend afin de mieux valoriser la contribution de recherche.
\end{itemize}

\section{Références Bibliographiques}
\nocite{*}
\bibliographystyle{ieeetr}
\bibliography{references}

\end{document}


